% Generated draft: requires mathematical and source review.
\documentclass[11pt]{article}
\usepackage[margin=25mm]{geometry}
\usepackage{fontspec}
\setmainfont{Latin Modern Roman}
\usepackage{amsmath,amssymb,mathtools,unicode-math}
\setmathfont{Latin Modern Math}
\usepackage{enumitem,needspace,xurl,hyperref}
\setlength{\parindent}{0pt}
\setlength{\parskip}{5pt}
\setlength{\emergencystretch}{3em}
\newcommand{\N}{\mathbb{N}}
\newcommand{\Z}{\mathbb{Z}}
\newcommand{\Q}{\mathbb{Q}}
\newcommand{\R}{\mathbb{R}}
\newcommand{\C}{\mathbb{C}}
\newcommand{\F}{\mathbb{F}}
\newcommand{\A}{\mathbb{A}}
\newcommand{\PP}{\mathbb{P}}
\newcommand{\E}{\mathbb{E}}
\newcommand{\eps}{\varepsilon}
\newcommand{\dd}{\,\mathrm{d}}
\begin{document}
\section*{UnsolvedMath ID: 700454}
\textbf{Source problem code:} catalogue entry
\subsection{Definitions and mathematical statement}
A group \(P\) is polycyclic if it has a finite subnormal series
\[
1=P_0\triangleleft P_1\triangleleft\cdots\triangleleft P_r=P
\]
whose factors \(P_i/P_{i-1}\) are cyclic. A group is linear if it admits an injective homomorphism into \(\operatorname{GL}_n(K)\) for some field \(K\) and some integer \(n\geq 1\).

Let \(A\) and \(B\) be arbitrary polycyclic groups, let \(C\) be a group, and let
\[
\iota_A:C\hookrightarrow A,\qquad \iota_B:C\hookrightarrow B
\]
be injective homomorphisms such that \(\iota_A(C)\triangleleft A\) and \(\iota_B(C)\triangleleft B\). Define their amalgamated free product by
\[
G=A\mathbin{*_{C}}B
=(A*B)\big/\big\langle\!\big\langle
\iota_A(c)\iota_B(c)^{-1}:c\in C
\big\rangle\!\big\rangle,
\]
where the double brackets denote normal closure in \(A*B\).

\textbf{Question (B. A. F. Wehrfritz).} For every such choice of \(A,B,C,\iota_A,\iota_B\), do there exist a field \(K\), an integer \(n\geq 1\), and an injective homomorphism
\[
\rho:G\hookrightarrow\operatorname{GL}_n(K)?
\]
Here the field and dimension may depend on the amalgam. The source specifies no characteristic restriction; the formulation uses the usual meaning of linearity over some field.

\textbf{Supplied qualifications.} The source reports that \(G\) is residually finite, meaning that for every \(g\neq 1\) there is a homomorphism \(q:G\to F\) to a finite group with \(q(g)\neq 1\), and attributes this assertion to G. Baumslag (1963). It also reports affirmative special cases using the wording that the amalgamated part is ``torsion-free, abelian'' (B. A. F. Wehrfritz, 1973) ``and nilpotent'' (M. Shirvani, 1981, unpublished). The punctuation and scope of ``torsion-free'' do not unambiguously determine the hypotheses of these special cases. Their wording is retained without asserting a resolved interpretation. These are supplied background assertions, not additional hypotheses of the general question or independently verified results.
\subsection{Short English statement}
Must every free product of two polycyclic groups, amalgamating subgroups that are normal in both factors, have a faithful finite-dimensional representation over some field? The supplied record labels this question open and reports residual finiteness and affirmative special cases. The precise torsion assumptions in the reported special cases require source review.
\subsection{Sources}
\begin{itemize}
\item Local UnsolvedMath registry, permanent ID \(700454\), exact source code \texttt{catalogue entry}, title \emph{Kourovka Notebook Problem 8.2}. The supplied record labels the problem \emph{open}; that label is source metadata, not an independently verified assessment.
\item \emph{The Kourovka Notebook (Unsolved Problems in Group Theory)}, Problem 8.2, proposed by B. A. F. Wehrfritz in the 8th issue (1982). The supplied background identifies the 21st issue edition (2026), \texttt{arXiv:1401.0300v48}. The supplied collection description names E. I. Khukhro and V. D. Mazurov as editors.
\item G. Baumslag, \emph{Trans. Amer. Math. Soc.} \textbf{106}, no. 2 (1963), 193--209. Supplied attribution for residual finiteness; no article title is supplied.
\item B. A. F. Wehrfritz, \emph{Proc. London Math. Soc.} \textbf{27}, no. 3 (1973), 402--424. Supplied attribution for the affirmative case described as ``torsion-free, abelian''; no article title is supplied.
\item M. Shirvani (1981), unpublished. Supplied attribution for the affirmative case described as ``nilpotent''; no title or further bibliographic information is supplied.
\end{itemize}
\end{document}
